\section{Invariants, Conservation Laws, and Property-Based Testing}
\label{sec:sec15}
\label{sec:invariants}% canonical home of the invariant prose; other sections \ref this id
%==============================================================================

This section is the consolidation point for the framework's correctness claims. Properties argued elsewhere through worked examples are restated here once, as precise testable invariants, in one reconciled numbering. They serve two roles at once: the design specification---the system is correct exactly when they all hold---and the test oracle against which property-based testing checks every generated input. Executable pre/postcondition forms attach in Appendix~\ref{app:properties}; this section states the prose and the cross-links.

%------------------------------------------------------------------------------
\subsection{Conservation as Central Invariant and Universal Oracle}
\label{sec:sec15-conservation}

The conservation law $Q(u)=\sum_w w_t(u)=0$, for every unit $u$ and time $t$, is the central invariant. It is an algebraic consequence of move semantics ($\mathrm{src}\mathrel{-}=q;\ \mathrm{dst}\mathrel{+}=q$) and holds for every transaction by construction (Section~\ref{sec:ledger}); the conserved-charge analogy is exact, the charge being the total quantity of each unit.

Conservation governs \emph{quantities per unit}, not values. An FX trade conserves EUR and USD quantities separately; at any rate other than the trade rate the value of EUR received differs from that of USD paid. Value changes from price movements, including FX, are carried by the valuation projection $V_t=\sum_u w_t(u)\cdot P_t(u)$, specified in the valuation/market-data spec (Section~\ref{sec:valuation}), not by conservation. Valuation is a linear functional on the state; dual valuation applies two price vectors to one state; PnL is the functional's difference at two times---all derived from conservation and the additive structure of moves.

Conservation is therefore a universal oracle: for any transaction sequence, checking $Q(u)=0$ is constant-time after the sum and detects any violation of the system's fundamental guarantee. A conservation failure is an implementation bug, never a financial event.

%------------------------------------------------------------------------------
\subsection{The Core Invariants P1--P10}
\label{sec:sec15-core}

The ten core invariants are the acceptance criteria for any implementation. An implementation that violates any one is incorrect regardless of its outputs. Appendix~\ref{app:properties} gives each as an executable oracle; the prose statement is fixed here.

\begin{enumerate}
\item[\textbf{P1}] \textbf{Conservation.} $Q(u)=\sum_w w_t(u)=0$ for all units $u$, all times $t$. (Stated as System Closure in Section~\ref{sec:ledger}, Theorem~\ref{inv:P1}.)
\item[\textbf{P2}] \textbf{Atomic commitment.}\label{inv:P2} A transaction applies completely or not at all; no partial state is observable.
\item[\textbf{P3}] \textbf{Referential integrity.}\label{inv:P3} Every move references existing wallets and an existing unit.
\item[\textbf{P4}] \textbf{Log monotonicity.}\label{inv:P4} The event stream is append-only; no entry is retroactively rewritten. Tamper-evidence is provided by hash chaining (each entry carries the hash of its predecessor), which detects tampering but does not prevent it; prevention requires deployment controls (WORM storage, consensus, or hardware security modules). The design requires tamper-evidence as the minimum; stronger immutability depends on the environment.
\item[\textbf{P5}] \textbf{Transaction idempotency.}\label{inv:P5} Applying the same transaction twice (by \texttt{tx\_id}) has no incremental effect.
\item[\textbf{P6}] \textbf{Lifecycle idempotency.}\label{inv:P6} Applying the same lifecycle event to the same unit state produces no additional moves.
\item[\textbf{P7}] \textbf{Virtual/real ledger isolation.}\label{inv:P7} No move has one endpoint in a virtual ledger and the other in the real ledger.
\item[\textbf{P8}] \textbf{Snapshot consistency.}\label{inv:P8} \texttt{clone\_at($t$)} reconstructs the complete view---balances \emph{and} unit state---identical to the ledger state at time $t$.
\item[\textbf{P9}] \textbf{Lifecycle purity.}\label{inv:P9} For fixed inputs, lifecycle functions produce identical outputs and no effect beyond the return.
\item[\textbf{P10}] \textbf{PnL path-independence.} $\mathrm{PnL}=V_{t_1}-V_{t_0}$, independent of the intermediate transactions (Section~\ref{sec:valuation}, Theorem~\ref{inv:P10}).
\end{enumerate}

Two further oracles the same generators exercise are global rather than numbered among the ten: \emph{totality} (every valid input returns a success or a well-typed error, never an unhandled exception) and \emph{valid transitions only} (every product/state/event triple absent from the product's transition table is rejected). Both are stated in executable form in Appendix~\ref{app:properties}.

%------------------------------------------------------------------------------
\subsection{SBL Invariants P11--P20}
\label{sec:sec15-sbl}

Securities borrowing and lending adds ten invariants that narrow the valid-state set without relaxing any core invariant: every transaction satisfying P11--P20 also satisfies P1--P10. They are stated in full, with the position vector and Single-Coordinate Move Principle they rest on, in Section~\ref{sec:sbl-invariants}; the consolidated index follows.

\begin{description}[leftmargin=3.2em,style=nextline]
\item[P11 Paired legs] Every SBL settlement transaction carries both a securities and a collateral leg (exceptions: uncollateralised, cash pool, triparty RQV).
\item[P12 On-loan consistency]\label{inv:P12} $w^{\mathrm{onloan}}_e(u)$ equals the summed quantity of that entity's live loans of $u$.
\item[P13 Collateral sufficiency]\label{inv:P13} Posted collateral covers exposure, by one of three formulations: per-loan cash, cash pool, non-cash bilateral.
\item[P14 Locate before short]\label{inv:P14} Every short sale satisfies a locate precondition.
\item[P15 Partial-return monotonicity]\label{inv:P15} Loan quantity is monotonically non-increasing.
\item[P16 SFTR completeness]\label{inv:P16} Each reportable SBL event produces exactly one SFTR report.
\item[P17 Settlement-state monotonicity]\label{inv:P17} Settlement states advance monotonically.
\item[P18 Lender ownership invariance]\label{inv:P18} A pure lender's $\mathrm{own}$ coordinate is constant across lending operations under title transfer (buy-in excluded). \emph{By construction:} under the Single-Coordinate Move Principle, loan initiation and return write $\mathrm{onloan}$, never the lender's $\mathrm{own}$.
\item[P19 Rehypothecation regime compliance]\label{inv:P19} Rehypothecation is validated against the legal regime.
\item[P20 Available-inventory identity]\label{inv:P20-hub} $\mathrm{avail}(e,u)=\vec w_e(u)[\mathrm{own}]-\vec w_e(u)[\mathrm{onloan}]+\vec w_e(u)[\mathrm{borr}]$. \emph{By construction:} a definition computed on read from three stored coordinates, never stored, hence not violable (Invariant~\ref{inv:P20}).
\end{description}

%------------------------------------------------------------------------------
\subsection{Obligation Liveness Invariants P21--P23}
\label{sec:sec15-obligation}

Event-triggered obligations\allowbreak---\allowbreak recalls, collateral substitutions, margin deliveries, close-out deadlines\allowbreak---\allowbreak are first-class objects governed by three invariants, stated in full in Section~\ref{sec:obligation-invariants}.

\begin{description}[leftmargin=3.2em,style=nextline]
\item[P21 Obligation liveness] For every obligation $o$ with deadline $t_d$: $\forall t>t_d,\ o.\mathit{state}\in\{\textsc{Discharged},\allowbreak\textsc{Compensated},\allowbreak\textsc{Defaulted}\}$. No obligation persists in \textsc{Pending} or \textsc{Attempted} past its deadline (Invariant~\ref{inv:P21}).
\item[P22 Obligation conservation] Every discharging and every compensating move set satisfies $Q(u)=0$; the mechanism does not bypass the executor. Follows from P1 (Invariant~\ref{inv:P22}).
\item[P23 Obligation idempotency] Discharging the same obligation twice (by $o.\mathit{id}$) has no incremental effect; the discharge predicate is checked before any move is emitted. Composes with P5 and P6 (Invariant~\ref{inv:P23}).
\end{description}

%------------------------------------------------------------------------------
\subsection{The Twelve Conditions C1--C12}
\label{sec:sec15-conditions}

The three-home state model (Section~\ref{sec:sec04}) imposes twelve conditions, each defined once at the instrument that forces it. Each renders an illegal state unrepresentable rather than merely tested. The index below cross-links them; the discharge map (Section~\ref{sec:sec15-discharge}) gives the mechanism by which each acts.

\begin{center}
\small
\begin{longtable}{@{}p{1cm}p{10.4cm}@{}}
\toprule
\textbf{Cond.} & \textbf{Content (defined in Section~\ref{sec:sec04}, except C7/C10 in Section~\ref{sec:unit-store})} \\
\midrule
\endhead
C1 & Position accessor returns Option; monotone carrier keeps the key set stable. \\
C2 & Handler-level conservation per event class, vacuous base case included. \\
C3 & Atomic \texttt{Transaction} across all three maps; partial application unrepresentable. \\
C4 & Capability-scoped reads; cross-overlay reads forbidden; exports via \emph{UnitStatus}. \\
C5 & \emph{UnitStatus} registration-total; product-declared defaults at registration. \\
C6 & \emph{ProductTerms} versioned, append-only; no in-place mutation. \\
C7 & \emph{ProductTerms} registration-total; first write at registration. \\
C8 & Two-track amendment by a total fungibility predicate (Preserving / Breaking). \\
C9 & Zero-holder handlers discharge $\sum=0$ vacuously. \\
C10 & Re-registration is a hard error, never a silent reset. \\
C11 & Each \emph{PositionState} field has a canonical writer set; an outside writer is a type error at authorship. \\
C12 & All per-(wallet, mandate/strategy) economic state lives in \emph{PositionState}; the $W$-sector is empty by schema. \\
\bottomrule
\end{longtable}
\end{center}

%------------------------------------------------------------------------------
\subsection{Which Condition Discharges Which Invariant}
\label{sec:sec15-discharge}

The highest-value content of the consolidation is the map from conditions to the guarantees they place beyond reach, and the mechanism in each case. The guarantees carry the canonical P1--P23 names this section fixes; Section~\ref{sec:sec04} presents the same map under its own local labels, for which the P1--P23 names here are canonical. ``Unrepresentable'' is used precisely: an \emph{abstract type} whose only constructor is a checked smart constructor makes an illegal value impossible to build, so it never reaches the operation that would act on it; a \emph{non-empty list type} makes a versionless state untypable; a \emph{GADT write-tag} is a per-field marker checked at compile time, then erased once writes share a stored row; a \emph{value-level check} is a runtime guard returning an error, not a type fact. Conservation, by contrast, holds by construction: each \texttt{Move} edge names both parties, so the signed per-unit sum of any move list is the group identity \texttt{mempty}---a structural fact of the \texttt{Transaction} type, not a runtime check.

\begin{center}
\small
\begin{longtable}{@{}p{3.5cm}p{1.6cm}p{6.4cm}@{}}
\toprule
\textbf{Guarantee} & \textbf{Cond.} & \textbf{Mechanism} \\
\midrule
\endhead
P1 Conservation & C2, C3 & A \texttt{Transaction} carries its flow as \texttt{Move} edges; each edge contributes $\mathrm{qneg}\,q\mathbin{<\!\!>}q=\texttt{mempty}$, so the signed per-unit sum of any move list is the group identity by construction---the empty list (registration) included. No unconserved \texttt{Transaction} exists to reject, hence no validate gate. C3 makes one \texttt{Transaction} span all three maps, so partial application is unrepresentable. \\
P8 Snapshot / replay & C1 & Replay is the \texttt{foldM} of the event stream---a catamorphism, meaning replaying the events always rebuilds the exact value---hence a homomorphism $\texttt{replay}(xs\mathbin{<\!\!>}ys)=\texttt{replay}\,xs \mathrel{>\!\!=\!\!>} \texttt{replay}\,ys$, where $\mathbin{>\!\!=\!\!>}$ runs the first error-returning step and feeds its result to the next, stopping at the first error. The monotone carrier (C1) keeps the key set stable across any checkpoint cut, so checkpoint-independence follows from the law, not a test. \\
P6 Lifecycle idempotency & C5 & The lifecycle stage lives in \emph{UnitStatus}$[u]$, $u$-keyed and written by replacement, so re-applying a stage write is the identity (\texttt{EXPIRED} over \texttt{EXPIRED} is \texttt{EXPIRED})---structural at one key. Non-conserved \emph{PositionState} fields share the algebra (\texttt{hwm} by $\max$, \texttt{entry\_nav} write-once); conserved fields draw replay-safety from P1. \\
Terms immutability (underpins P4) & C6, C7 & \emph{ProductTerms} is an append-only $\mathrm{NonEmptyList}[\mathit{TermsVersion}]$, so a versionless unit is untypable; registration totality fixes the first version at registration. \\
Identity immutability (underpins P4) & C10, C8 & Re-registration is a hard error; a breaking amendment allocates a new unit $u_{\mathrm{new}}$ and never rewrites $u_{\mathrm{old}}$. \\
P7 Isolation / read scoping & C4 & Isolation holds by the per-move locality predicate (both endpoints in one ledger), not by conservation; cross-overlay reads are forbidden at the capability layer wrapping the accessors, not in the data reference. \\
Handler-field canon (underpins P9) & C11 & The canonical writer set per field is a \texttt{FieldWrite} GADT: a write by a field-writer outside the set is a type error at authorship, erased once writes share a delta row. \\
\bottomrule
\end{longtable}
\end{center}

The lifecycle-idempotency and snapshot guarantees rest on the status discipline:

\begin{remark}[UnitStatus is a projection, not a source of truth]
UnitStatus holds one value per unit, shared---read identically by every holder. Its value changes over time, but UnitStatus is not a separate source of truth: the immutable event log is. Every change is caused by a logged event; the stored value is overwritten in place only as a read cache. Replaying the events up to any point rebuilds the exact value that held then, so nothing is lost by overwriting and there is no other way to change the value. A value exists from registration onward.
\end{remark}

Each mechanism above is a type or a value-level gate carrying a theorem; the Haskell rendering makes the correspondence explicit.

\noindent\textbf{Types as theorems.} The five anchors below render the discharge mechanisms as Haskell, one type or function per guarantee. They are verbatim excerpts of the FORMALIS-cleared reference (\texttt{reference/Ledger.hs}, the three-home model of Section~\ref{sec:sec04}; see Appendix~\ref{sec:appB}); each carries that clearance. Quantities are exact \texttt{Integer} minor units forming an additive abelian group (\texttt{Qty}, \texttt{<>}~=~addition, \texttt{mempty}~=~zero); floating point is excluded because conservation and replay are arithmetic facts it would forfeit.

\paragraph{(1) P1 conservation --- \texttt{Move} edges / \texttt{unitDelta} (C2).} Conservation is a property of the \texttt{Transaction} type, not a checked side-condition. Each \texttt{Move} names both parties, so its per-unit contribution is $\mathrm{qneg}\,q\mathbin{<\!\!>}q=\texttt{mempty}$; summing over any move list is a \emph{group homomorphism}---the total over all wallets equals the running sum taken edge by edge---landing at the group identity \texttt{mempty} by the monoid laws. The empty list (a move-less registration or amendment) is \texttt{mempty}, so a zero-holder event (C9) conserves with no special case; summing edges, never dividing by a holder count, excludes the divide-by-zero-holders bug class. There is no unconserved \texttt{Transaction} to reject, hence no validate gate.

\begin{lstlisting}[language=Haskell]
data Move = Move
  { mFrom :: !WalletId, mTo :: !WalletId, mUnit :: !UnitId
  , mQty  :: !Qty, ... }                  -- positive magnitude; from -= q, to += q

-- Conservation witness: the signed per-unit sum over the edges. By construction:
unitDelta :: UnitId -> Transaction -> Qty
unitDelta u tx =
  foldMap (\m -> if mUnit m == u then qneg (mQty m) <> mQty m else mempty) (txMoves tx)
-- law:  unitDelta u tau == mempty   for every unit u and transaction tau

netDelta :: Transaction -> Qty
netDelta tx = foldMap (\m -> qneg (mQty m) <> mQty m) (txMoves tx)
-- law:  netDelta tau == mempty   for every tau (the empty/registration case too)
\end{lstlisting}

\noindent A one-edge transfer shows there is nothing to reject: a \texttt{Move} names both endpoints, so a half-edge is unrepresentable.
\begin{lstlisting}[language=Haskell]
-- tx = Transaction u [Move a b u 1000 t s] mempty [] Nothing Nothing
netDelta tx == mempty           -- by construction; a one-sided delta cannot be built
\end{lstlisting}

\paragraph{(2) P4 terms immutability --- \texttt{NonEmpty} (C6/C7).} \texttt{ProductTerms} wraps a \emph{non-empty} list, so a ``registered but versionless'' unit is untypable; \texttt{currentTerms} is total without a \texttt{Maybe}. No setter exists; the only growth operations are \texttt{register} (singleton) and \texttt{appendVersion} (append, never rewrite).

\begin{lstlisting}[language=Haskell]
newtype ProductTerms = ProductTerms (NonEmpty TermsVersion)    -- abstract

currentTerms :: ProductTerms -> TermsVersion                  -- total: no Maybe
currentTerms (ProductTerms vs) = NE.last vs

appendVersion :: TermsVersion -> ProductTerms -> ProductTerms -- append-only (C6)
appendVersion tv (ProductTerms vs) = ProductTerms (vs <> (tv :| []))
\end{lstlisting}

\paragraph{(3) P8 monotone carrier --- the Option accessor (C1).} \texttt{position} returns \texttt{Maybe}: \texttt{Nothing} is \emph{never held}, \texttt{Just zeroP} is \emph{held and currently flat}---distinct states the spec reads apart. The carrier is monotone by absence: \texttt{Ledger} is abstract, \texttt{applyTx} only inserts or updates a row, and no row-deleter is exported, so the key set never shrinks and a checkpoint cut cannot lose a row.

\begin{lstlisting}[language=Haskell]
position :: Ledger -> WalletId -> UnitId -> Maybe PositionState
position l w u = Map.lookup (w, u) (ledgerPS l)
--  Nothing    = wallet has NEVER held this unit
--  Just zeroP = held once, currently flat   (the two cannot be collapsed)
\end{lstlisting}

\paragraph{(4) P8 replay homomorphism --- Kleisli \texttt{>=>}.} Replay is an error-stopping fold (a \emph{catamorphism}---meaning replaying the events always rebuilds the exact value). Determinism and checkpoint-independence are the homomorphism law $\texttt{replay}(xs\mathbin{<>}ys)=\texttt{replay}\,xs\mathrel{>\!\!=\!\!>}\texttt{replay}\,ys$, not a test: replaying two batches in sequence equals replaying them together. With the stable key set from (3), where a cut falls cannot change the result.

\begin{lstlisting}[language=Haskell]
-- (>=>) runs the first Either-returning step, feeds its result to the next,
-- and stops at the first error.
replay :: [Transaction] -> Ledger -> Either LedgerError Ledger
replay ds l0 = foldM (\l d -> applyTx d l) l0 ds
\end{lstlisting}

\paragraph{(5) Handler-field canon --- \texttt{FieldWrite} GADT (C11, underpins P9).} The GADT constructors \emph{are} the field-to-writer table: a write by a handler outside a field's writer set is a type error at authorship. A handler is typed by the \texttt{Handler} it speaks for, so the phantom index constrains at a live call site, not merely in an alias; the index is erased to \texttt{SomeWrite} once authored, the check having already happened.

\begin{lstlisting}[language=Haskell]
data Handler = Settle | Trade | Transfer | FeeCrystallise | Subscribe

-- The balance field is NOT in this table: its sole canonical writer is the Move
-- edge (handler 'Transfer), applied in applyTx -- so a balance change without a
-- conserving move is unrepresentable. The rest are non-balance per-position rows.
data FieldWrite (h :: Handler) where
  WAc       :: Qty -> FieldWrite 'Settle
  WAcTrade  :: Qty -> FieldWrite 'Trade
  WHwm      :: Qty -> FieldWrite 'FeeCrystallise
  WEntryNav :: Qty -> FieldWrite 'Subscribe

settleHandler :: [(WalletId, Qty)] -> Map WalletId [FieldWrite 'Settle]
settleHandler legs = Map.fromList [ (w, [WAc q]) | (w, q) <- legs ]
-- A settle body bumping hwm would emit WHwm :: FieldWrite 'FeeCrystallise and
-- fail to typecheck against the declared 'Settle output -- C11 IS a type error.
\end{lstlisting}

%------------------------------------------------------------------------------
\subsection{Invariants as Specification and Oracle}
\label{sec:sec15-spec}

The invariants define correctness independently of any implementation: the system is correct if and only if all hold, and any property-based test checking an invariant thereby tests the specification. An example-based test verifies one input/output pair; an invariant-based test verifies a property over a whole input class.

A bond coupon payment illustrates the dual role. P1 requires $Q(\mathrm{USD})=0$ before and after the coupon move; P6 requires that paying the same coupon twice produces no additional moves. Applied to randomly generated bonds with random coupon schedules, these two checks exercise scenarios no hand-written suite would enumerate. The equation $Q(u)=0$ governs both the ledger and its test suite.

%------------------------------------------------------------------------------
\subsection{Property-Based Testing over the CDM Generator Universe}
\label{sec:sec15-pbt}

Each property has a \emph{generator} (the random-input space) and a \emph{postcondition} (the invariant). The framework draws inputs, applies the operation, and checks the postcondition; failure yields a minimal counterexample. The executable pre/postcondition for each of P1--P10, together with the totality and valid-transition oracles, is given once in Appendix~\ref{app:properties}.

The generators draw from the CDM enumerations---product types, event types, party roles, date rules (Section~\ref{sec:cdm}). Because these enumerations are closed and finite, the state space of any product's lifecycle is bounded, and coverage is a checklist rather than a probability: where a random fuzzer might miss a rare combination after millions of runs, the enum-driven generator covers every combination in its first pass. The strategy is compositional: draw a CDM product, draw a sequence of valid lifecycle events, map through $F$ to ledger transactions, and check every invariant.

\paragraph{Completeness by enumeration.} The CDM \texttt{EventIntentEnum} is the closed set of lifecycle intents. For an equity option the applicable intents map one-to-one to unit-state transitions:

\begin{center}
\begin{tabular}{lll}
\toprule
\textbf{EventIntentEnum} & \textbf{Unit-state transition} & \textbf{Status} \\
\midrule
\texttt{OPEN} & $\emptyset\to$ ACTIVE & Covered \\
\texttt{EXERCISE} & ACTIVE $\to$ EXERCISED & Covered \\
\texttt{EXPIRE} & ACTIVE $\to$ EXPIRED & Covered \\
\texttt{ASSIGN} & ACTIVE $\to$ ASSIGNED & Covered \\
\texttt{NOVATION} & ACTIVE $\to$ NOVATED & Covered \\
\texttt{PARTIAL\_TERMINATION} & --- & Not applicable (option is all-or-nothing) \\
\bottomrule
\end{tabular}
\end{center}

Because the enum is closed, the generator drawing intents uniformly from it achieves full intent coverage by construction. Incompleteness is surfaced, not hidden: a new enum value with no transition (say \texttt{RESTRUCTURING}) is produced by the generator, finds no matching rule, and is flagged by the valid-transitions oracle. A missing rejection is as much a bug as a missing acceptance. The \texttt{OptionTypeEnum} $\{$\texttt{CALL}, \texttt{PUT}, \texttt{PAYER}, \texttt{RECEIVER}, \texttt{STRADDLE}$\}$ bounds payout directionality the same way: every variant must have payout logic, and a missing case is a compile-time or test-time error, not a runtime surprise.

%------------------------------------------------------------------------------
\subsection{Specification-First Development}
\label{sec:sec15-specfirst}

The properties are written before the implementation: typed state schemas make illegal field combinations unrepresentable; the core properties are executable specifications; the simplest lifecycle case anchors the simulation harness; the full product universe is reached using CDM enumerations as the generator source. Correctness is thereby defined independently of language, storage backend, or deployment, and the same properties hold from single-process prototype to distributed system. The approach is falsifiable: a property satisfiable by no reasonable implementation reveals a design contradiction before implementation begins.

%==============================================================================
